Die vorstehenden Ausführungen sollten die lebenswichtige Bedeutung der Progressionsregelung bei einer Gesellschaft mit erheblich wachsender Teilnehmerzahl und damit die Wahrheit meiner auf S. [127] aufgestellten Behauptung darlegen. Um diesen Gegenstand jedoch von allen Seiten zu beleuchten, muss der Entstehung des betreffenden Statutsartikels Schritt für Schritt nachgegangen werden. Ich habe dazu die betreffenden Akten sorgfältig gelesen und wiederholt gelesen und gebe daraus einen Auszug, soweit sie jenen Statutsartikel betreffen. Dabei werde ich hier und da auch einige an sich untergeordnete Nebenumstände hervorheben müssen, wenn sie dazu beitragen können, den Verlauf dieser Verhandlungen verständlicher zu machen. Vorab bemerke ich, dass 1794 der jährliche Beitrag 5 Thl. Gold betrug, die Witwenpension 110 Thl. Kassenmünze und dass vater- und mutterlose Waisen bis zum vollendeten \(12^{\text{ten}}\) Jahr Anspruch auf die Pension hatten.

In einem an die Universität gerichteten Ministerialreskript vom 30. Juni 1794 wurde unter Hinweis auf ein schon einige Zeit zuvor vom König an die Witwenkasse gegebenes Geschenk von 1000 Thl.\footnote{Kassenmünze. Nach der Rechnung für 1793 wurde es am 18. April 1793 eingezahlt.} die Bewilligung eines zweiten Geschenks von 500 Thl. Gold mitgeteilt. Dabei wurde hinzugefügt, dass diese Geschenke eine baldige Pensionserhöhung ermöglichen sollten und man deshalb erwarte, dass die Teilnehmer auch ihrerseits zu diesem Zweck beitragen und bereit sein würden, die jährlichen Beiträge von 5 Thl. Gold auf 10 Thl. Kassenmünze zu erhöhen. Nach den Berechnungen in einem beiliegenden P. M. sei es nicht zweifelhaft, dass schon jetzt ohne Weiteres eine Pensionserhöhung so eingeführt werden könne, dass die sechs ältesten Witwen statt der bisherigen 110 Thl. künftig 150 Thl. und alle übrigen jeweils 130 Thl. erhielten. Zum Schluss wurde angeboten, falls nicht sämtliche Kapitalien der Witwenkasse vollkommen sicher angelegt seien, ihre sichere Unterbringung zu 3 Prozent bei öffentlichen oder städtischen Kassen zu veranlassen. Ein ähnliches Angebot war bereits einmal bei den Bemerkungen zur Jahresrechnung für 1792 gemacht worden.

Die beigefügte Anlage, deren Verfasser nicht genannt ist, im Einzelnen durchzugehen, ist für meinen Zweck nicht nötig. Ein paar Nebenumstände will ich jedoch hervorheben.

I. Die Kapitalien der Witwenkasse seien, heißt es, zu unterschiedlichen Zinsfüßen ausgeliehen, einige\footnote{Damals beinahe der dritte Teil des Kapitalvermögens.} zu 3 Prozent, andere höher. Da der Zinsfuß für sämtliche Kapitalien aber leicht auf 3 Prozent zurückgehen könne und bei Überschlagsrechnungen mit möglichst sicheren Summen gerechnet werden müsse, wolle man bei den Berechnungen auch nicht mehr als 3 Prozent voraussetzen.

Daraus und aus dem eben angeführten Schluss des Reskripts erklärt sich, warum auch in den Verhandlungen bei der Universität für jeden künftigen Kapitalzuwachs ohne weitere Bemerkung nur mit 3 Prozent gerechnet wurde. Nur für die schon vorhandenen und bereits angelegten Kapitalien wurden die Zinsen mit \(3\tfrac12\) Prozent angesetzt. Vergleiche dazu die Anmerkung zu S. [131.]

II. Um zu überschlagen, auf welche Witwenzahl die Rechnung gestellt werden müsse, fährt der Verfasser fort:

‘Will man nun nach den bisherigen Erfahrungen annehmen, dass 3 bestehende Ehen eine Witwe zu er-\par
\noindent ‘nähren haben, so würden die 26 verheirateten Professoren (unter einer Gesamtzahl von 36) etwa 9 Witwen\par
\noindent ‘zu unterhalten haben. Da es aber größere Sicherheit bietet, den äußersten Fall zugrunde zu legen,\par
\noindent ‘setze man lieber für \(2\tfrac12\) Ehen \textit{eine} Witwe an, sodass die bestehenden\par
\noindent ‘26 Professorenehen — 10 Witwen zu unterhalten hätten.’
\par\SourceLeaf{138}
Bei den vonseiten der Universität gemachten Überschlagsrechnungen hielt man jene 3 oder \(2\tfrac12\) Ehen auf eine Witwe für zu viel und legte das Verhältnis von 2 Ehen auf eine Witwe zugrunde.

Ich habe die Stelle des P. M. hier nur deshalb angeführt, weil dadurch verständlich wird, dass man sich mit dem Verhältnis von zwei Ehen auf eine Witwe so leicht zufriedengab. Sehr wahrscheinlich entspricht aber auch dieses den Verhältnissen der Professorenwitwenkasse noch nicht; vielmehr sollten noch weniger Ehen auf eine Witwe gerechnet werden. Näheres zur Sache selbst wird weiter unten folgen. Der Verlauf des Geschäfts erinnert jedoch (wenn es erlaubt ist, hier einen Vergleich aus einem niedrigeren Bereich heranzuziehen) unwillkürlich an Käufer, die bei unvollkommener eigener Warenkenntnis glauben, ein gutes Geschäft gemacht zu haben, wenn sie weit unter dem zuerst verlangten Preis eingekauft haben, obwohl sie bei näherem Hinsehen noch immer zu teuer bezahlt haben. Ich brauche nicht daran zu erinnern, dass ich diesen Vergleich nicht über Gebühr ausgedehnt wissen will; denn der unbekannte Antragsteller hat die Verhältnisse \(3:1\) und \(2\tfrac12:1\) ohne Zweifel in gutem Glauben an ihre Zulässigkeit vorgebracht.

Als der damalige Prorektor F. am 6. Juli das Reskript im Senat in Umlauf setzte, fügte er den beiden darin enthaltenen Beratungsgegenständen (Erhöhung der Beiträge und Erhöhung der Pensionen) noch einen dritten hinzu: den Vorschlag, die Waisenpensionen bis zum vollendeten \(20^{\text{sten}}\) Lebensjahr zu verlängern. Er überließ es den Senatsmitgliedern, sich dazu gleich schriftlich oder in der für den 12. Juli angesetzten Senatsversammlung zu äußern.

Dieses Rundschreiben ist von 17 Senatsmitgliedern unterzeichnet; von den dabei abgegebenen Äußerungen sind hier nur wenige zu erwähnen.

Der damalige Kurator der Witwenkasse, P., stellt vor allem den Grundsatz auf, die Kasse sei den gegenwärtigen Witwen ebenso viel schuldig wie den künftigen, sie sei aber auch den künftigen Witwen genau so viel schuldig wie den gegenwärtigen. Diese (an sich tatsächlich sehr vage) Formulierung erläutert er so, dass jede künftige Witwe, die weniger erhalte als eine andere früher erhalten habe, (seiner Meinung nach) wirklich geschädigt werde. Der erste Grundsatz müsse daher sein, die Pensionshöhe so zu bestimmen, dass sie nach höchster Wahrscheinlichkeit niemals wieder vermindert zu werden brauche. In dieser Hinsicht hält P. die Berechnung in der Anlage des Reskripts aber nicht für sicher genug. Man dürfe nicht \(2\tfrac12\) Ehen auf 1 Witwe rechnen, sondern nur 2, müsse also das Maximum der Witwen nicht auf 10—12, sondern auf 14—15 festsetzen und folglich auch geringere Pensionshöhen annehmen.

\textsc{Kaestner} hält die Frage für zu verwickelt und schwierig, als dass ohne eine \textit{ausführliche} und \textit{genaue} Untersuchung etwas festgesetzt werden könne. Auch er sei der Ansicht, dass nur \textit{höchstens} 2 Ehen auf 1 Witwe gerechnet werden dürften. Da er schon lange aus der Witwenkasse ausgeschieden sei (er war von 1755—1773 Teilnehmer gewesen), habe er in der Sache keine Stimme (als Senior der philosophischen Fakultät war er allerdings Mitglied der Kirchendeputation). Er rate aber, keinen Beschluss zu fassen, ohne vorher einen Sachverständigen, etwa den p. \textsc{Kritter}, zu befragen.

G. wünscht ebenfalls, die Mittel der Kasse durch genaue Berechnungen zu ermitteln, und weist anhand der hannoverschen, bremischen und anderer Witwengesellschaften auf die \textit{schrecklichen} Folgen einer übereilten Bewilligung zu großer Pensionen hin.

Die übrigen Voten stimmen teils den vorigen zu, teils entwickeln sie Bedenken wegen der Beitragserhöhung oder der Verlängerung der Waisenpensionen; dies braucht hier nicht auszugsweise wiedergegeben zu werden.

In der Senatssitzung vom 12. Juli, an der einschließlich des Prorektors 12 Professoren teilnahmen, erklärten sich 9 unbedingt \textit{für} die Erhöhung der Beiträge, 2 mit dem Zusatz, dass sie die Erhöhung für zu groß hielten. Einer, der sich bei der schriftlichen Abstimmung nachdrücklich dagegen ausgesprochen hatte, wollte sich der Mehrheit anschließen. — F.s Vorschlag zur Verlängerung der Waisenpensionen wurde bis zu einer genaueren Prüfung der Kassenverhältnisse noch zurückgestellt. — Wegen der
\SourceLeaf{139}
Erhöhung der Witwenpensionen beschloss man, ein Gutachten von \textsc{Kritter} einzuholen\footnote{Es scheint nicht, dass etwas darüber festgelegt worden wäre, \textit{in welchem Maß} \textsc{Kritter}s Rat in Anspruch genommen werden sollte. \textsc{Kaestner} war bei der Versammlung nicht anwesend.}. Doch sprach man sich bereits \textit{gegen} die im Reskript vorgeschlagene höhere Pension für die sechs ältesten Witwen aus, da diese schon durch V.s Vermächtnis bevorzugt seien.

Die Befragung des p. \textsc{Kritter}, die durch ein als Entwurf bei den Akten liegendes Schreiben P.s geschah, war tatsächlich nur sehr beschränkt. Nach einer bloß in ganz allgemeinen Umrissen gehaltenen Übersicht der Haupteinrichtungen der Witwenkasse werden \textsc{Kritter} nur zwei Fragen vorgelegt: I) Nach welchem Grundsatz und Verhältnis in einer solchen Gesellschaft das Maximum der Witwen bestimmt werden müsse, und II) um wie viel dieses Maximum der auszuzahlenden Witwenpensionen noch vergrößert werden müsse, wenn bei Nichtvorhandensein einer Witwe oder nach deren früherem Tod der Pensionsanspruch auf etwa vorhandene Waisen übergehe und fortdauere, bis das jüngste Kind das Alter von 20 Jahren erreicht habe. Die damalige Gesamtzahl der Beteiligten und die Zahl der Verheirateten unter ihnen werden genau wie in der oben angeführten Anlage des Reskripts mit 36 und 26 angegeben. Zugleich wird bemerkt, dass diese Zahlen und ihr Verhältnis veränderlich und schwerlich einer Wahrscheinlichkeitsregel zu unterwerfen seien. Zum Schluss folgt das Angebot, weitere Daten aus den bisherigen Erfahrungen über das Verhältnis der Teilnehmer zu den Witwen sogleich mitzuteilen, falls der Befragte sie noch benötigen sollte.

\textsc{Kritter}s Antwort, oder sein »Gutachten«, vom 19. Juli lege ich in vollständiger Abschrift bei\footnote{Gutachten über einige mir vorgelegte Fragen zur Göttinger Universitätswitwenkasse.

Ad I. Das Kollegium der Herren Professoren in Göttingen besteht seit der Errichtung der Universität schon über 50 Jahre, sodass bereits vor 10 Jahren die höchste Witwenzahl vorhanden sein konnte, die nach den Gesetzen der Sterblichkeit auf etwa 22 oder 24 verheiratete Professoren kommen musste, nämlich 12 Witwen. Dies hat sich auch nach der mir mitgeteilten Liste über den Anwachs der jährlich vermehrten Witwenzahl gezeigt und hätte sich noch besser gezeigt, wenn das Kollegium der Herren Professoren aus 100 Personen hätte bestehen können. Da diese Zahl aber der Natur der Sache nach nur ungefähr den \(4^{\text{ten}}\) Teil von 100 ausmachte, war es auch natürlich, dass die Witwenzahl von 1775 bis 1780 mehr als 12 und von 1781 bis 1792 weniger als 12 betrug, weil bei kleinen Zahlen die Ordnung der Sterblichkeit nicht so eintreten kann wie bei großen. Dennoch muss man annehmen, dass die höchste Witwenzahl im Durchschnitt von etwa 10 Jahren beständig etwa halb so groß sein werde wie die Zahl der verheirateten Herren Professoren.

Ad II. Da nun gewünscht wird, dass die hinterlassenen Waisenfamilien einer verstorbenen Witwe oder auch eines verstorbenen Witwers bis zum vollendeten \(20^{\text{sten}}\) Jahr eine gleiche Pension wie eine Witwe erhalten möchten, kann ich darüber nur nach der Erfahrung bei der bremischen Witwenkasse etwas bestimmen. Die bremische Gesellschaft hatte ungefähr den \(6^{\text{ten}}\) Teil mehr an Witwenpensionen zu erwarten, als sie beschloss, Waisenfamilien wie eine Witwe zu pensionieren und damit bis zum vollendeten \(18^{\text{ten}}\) Jahr des jüngsten Kindes fortzufahren. Da die Waisenfamilien der Herren Professoren aber bis zum vollendeten \(20^{\text{sten}}\) Jahr des jüngsten Kindes in diesen Genuss kommen sollen, müssen doch wohl auf 12 Witwenpensionen mehr als 2 Waisenpensionen gerechnet werden.

Ich halte es deshalb für ratsam, zu 12 als dem Maximum der Witwenzahl noch wenigstens 2 hinzuzunehmen, sodass der Divisor, durch den die jährlichen Einkünfte aus den Fonds der Kasse und aus anderen Quellen geteilt werden, auf 14 gesetzt wird. Da jetzt nur 10 Witwen vorhanden sind, darf man die gesamten jährlichen Einnahmen der Kasse nicht unter diese 10 verteilen, weil die später hinzukommenden Witwen sonst benachteiligt würden; man muss vielmehr durch 14 teilen. Dann wird innerhalb von 10 Jahren der Fonds auf etwa 25 Witwenanteile verstärkt sein, deren Zinsen für eine Pension mehr als zuvor ausreichen werden, sodass man, wenn das wirkliche Maximum der Witwen- und Waisenpensionen eintritt, 15 Pensionen bezahlen kann. Sollte noch etwas übrig bleiben, könnte vor allem armen Witwen etwas hinzugezahlt werden.

Mögliche Unglücksfälle bei den angelegten Kapitalien, Zinsverluste und andere Ausgaben müssten doch wohl ebenfalls berücksichtigt werden.

Göttingen, den 19. Juli 1794.\hfill J. A. \textsc{Kritter}.}.

